\chapter{DESCONEXIÓN EMOCIONAL EN LAS REDES SOCIALES}
\section{Uso de redes sociales y soledad en jóvenes}

La relación entre el uso de redes sociales y la soledad ha sido objeto de numerosas investigaciones en los últimos años, con resultados que durante mucho tiempo resultaron contradictorios. Sin embargo, la evidencia más reciente permite precisar en qué condiciones esta relación se manifiesta.
Feng, Lien, Cian, Wang, Li, Chen y Wu (2026) realizaron un meta-análisis que sintetizó 59 estudios con 67 muestras independientes y un total de 81,344 participantes adolescentes y jóvenes adultos (edad media de 17.6 años), publicados entre el año 2000 y mayo de 2025. Los autores encontraron que el uso problemático de redes sociales ---entendido como un patrón de uso compulsivo y difícil de controlar--- mostró una asociación moderada y consistente con la soledad ($r = 0.29$), mientras que el tiempo total dedicado a las plataformas mostró una correlación considerablemente menor ($r = 0.07$). Este hallazgo es central para el presente proyecto, pues sugiere que la desconexión emocional juvenil no depende principalmente de la cantidad de horas de exposición a las redes sociales, sino del tipo de vínculo que la persona establece con su uso.
En la misma línea, Godard y Holtzman (2024) desarrollaron un meta-análisis de 141 estudios con aproximadamente 145,000 participantes, centrado en comparar los efectos del uso activo (crear contenido, comentar, interactuar directamente con otros) frente al uso pasivo (observar contenido sin interactuar). Sus resultados mostraron que el uso activo se asoció de manera significativa con un mayor apoyo social percibido en línea ($r = 0.34$), mayor bienestar general ($r = 0.15$) y mayor afecto positivo ($r = 0.11$), mientras que la mayoría de los tamaños de efecto para otras variables resultaron insignificantes. Este hallazgo complementa al anterior: si bien el tiempo de uso por sí solo no predice la soledad, el tipo de interacción ---particularmente la interacción activa y recíproca--- sí se relaciona con mejores resultados emocionales.


\section{Tipo de interacción digital como factor determinante}

Más allá de establecer que existe una relación entre uso de redes sociales y bienestar emocional, resulta necesario identificar con precisión qué características de la interacción digital explican dicha relación.
Liu, Pradhan y Wang (2026) realizaron una revisión sistemática siguiendo el protocolo PRISMA, en la cual analizaron 21 estudios empíricos publicados entre 2000 y 2024 sobre la asociación entre el uso de redes sociales y el aislamiento social en jóvenes adultos. Los autores organizaron los hallazgos en tres ejes temáticos: las motivaciones y patrones de uso de las plataformas, los efectos diferenciados de los canales de comunicación textual frente a los visuales, y la ausencia de asociación significativa en un subconjunto de estudios. Este último eje resulta relevante porque evidencia que la relación entre redes sociales y aislamiento no es uniforme, sino que depende de variables contextuales como el tipo de canal utilizado.
Un estudio exploratorio realizado por Tian (2026) con 150 participantes chinos, si bien de alcance metodológico más limitado por el tamaño y la homogeneidad de su muestra, aporta evidencia consistente con lo anterior en un contexto cultural distinto: mediante el \textit{Social Media Activity Questionnaire} (SMAQ), encontró que el uso activo se correlacionó negativamente con la comparación social, la depresión, la ansiedad y el estrés, mientras que el uso pasivo mostró el patrón inverso en las cuatro variables. Este resultado, aunque no generalizable por las limitaciones de la muestra, refuerza la idea de que el tipo de interacción ---y no solamente la exposición a la plataforma--- es un factor determinante del estado emocional del usuario.
En conjunto, esta evidencia permite identificar un primer factor relevante para el diseño de la aplicación: los mecanismos que fomentan una interacción activa, directa y recíproca entre los usuarios tienden a asociarse con mejores resultados emocionales que aquellos que promueven el consumo pasivo de contenido.


\section{Conexión abundante versus vínculo significativo}

Un segundo eje de análisis se refiere a la naturaleza cualitativa de los vínculos que las plataformas digitales permiten construir. Smith, Leonis y Anandavalli (2021) plantean que la soledad no debe entenderse como la ausencia de contactos, sino como la discrepancia percibida entre las relaciones sociales deseadas y las relaciones efectivamente disponibles. Bajo esta perspectiva, un joven puede acumular cientos de contactos digitales y experimentar, al mismo tiempo, un nivel elevado de soledad, en la medida en que esos contactos no satisfacen su necesidad de conexión genuina.
Esta idea encuentra respaldo conceptual en el modelo propuesto por Takano (2018), quien distingue entre dos formas de vinculación social: el cuidado social elaborado (\textit{elaborate social grooming}), dirigido a relaciones cercanas y que da lugar a redes sociales reducidas pero profundas, y el cuidado social ligero (\textit{lightweight social grooming}), característico de las redes sociales digitales, que permite mantener un número mucho mayor de contactos a costa de la profundidad del vínculo. Este modelo ofrece un marco explicativo para comprender por qué las plataformas actuales, orientadas a maximizar el número de conexiones, tienden a producir relaciones numerosas pero superficiales.
Evidencia cualitativa reciente ilustra cómo los propios usuarios experimentan esta tensión. Parry, Kuit, Murray y van der Westhuizen (2026) analizaron las narrativas de una comunidad en línea orientada a reducir el uso de tecnología digital, y encontraron que sus miembros describen la tecnología simultáneamente como una vía de escape y como una fuente de angustia, desarrollando distintas estrategias para gestionar su relación con ella. Un tema recurrente en estas narrativas fue lo que los autores denominan la ilusión de la amistad significativa en línea, es decir, la percepción de que los vínculos construidos exclusivamente en el entorno digital no logran replicar la profundidad de una relación presencial.


\section{Resultados de la encuesta diagnóstico aplicada al grupo objetivo}

\textit{Esta sección se completa una vez aplicada la encuesta de diagnóstico descrita en el apartado 1.6.2 del presente documento, dirigida a jóvenes de 16 a 28 años. Se recomienda estructurar la presentación de resultados en los siguientes puntos:}

\begin{itemize}
    \item Caracterización de la muestra (edad, género, plataformas de uso frecuente).
    \item Nivel percibido de desconexión emocional al interactuar con redes sociales actuales.
    \item Percepción sobre la calidad versus la cantidad de contactos digitales.
    \item Valoración de elementos como métricas públicas (likes, seguidores, vistas) en la experiencia emocional del usuario.
    \item Interés y expectativas frente a una herramienta centrada en micro interacciones breves y significativas.
\end{itemize}

Los resultados obtenidos deberán contrastarse con los hallazgos reportados en las secciones 3.1 a 3.3, de manera que el diagnóstico no se sustente únicamente en la literatura, sino también en evidencia primaria recogida en el contexto local del proyecto.

\section{Factores identificados que limitan la generación de vínculos significativos}

A partir de la evidencia revisada, es posible sintetizar los siguientes factores que contribuyen a la desconexión emocional de los jóvenes en su interacción con las plataformas digitales actuales:

\begin{enumerate}
    \item \textbf{Uso problemático más que uso extendido.} No es el tiempo de exposición a las redes sociales lo que se asocia con mayor soledad, sino un patrón de uso compulsivo y difícil de autorregular (Feng et al., 2026).
    \item \textbf{Predominio del consumo pasivo sobre la interacción activa.} El consumo pasivo de contenido ---observar sin interactuar--- se asocia consistentemente con peores resultados emocionales que la interacción activa y recíproca (Godard \& Holtzman, 2024; Tian, 2026).
    \item \textbf{Ausencia de mecanismos que prioricen la calidad del vínculo.} Las plataformas actuales están diseñadas para maximizar el número de conexiones y el tiempo de uso, sin incorporar mecanismos que favorezcan la profundidad de la interacción (Takano, 2018).
    \item \textbf{Presencia de métricas de comparación social.} Los elementos de validación pública propios de las redes sociales convencionales (likes, seguidores, vistas) están vinculados en la literatura a la comparación social ascendente, favoreciendo la sensación de inadecuación más que la conexión genuina.
    \item \textbf{Percepción de superficialidad en los vínculos digitales.} Los propios usuarios identifican una brecha entre la cantidad de contactos digitales disponibles y la calidad de conexión emocional que estos les proporcionan (Smith et al., 2021; Parry et al., 2026).
\end{enumerate}